According to variation theory \autocite[See][]{NCOL}, learning an educational
objective requires the learner to discern all critical aspects of that
objective. Misconceptions occur when the learner cannot yet discern one or
more critical aspects; documented misconceptions therefore help us identify
what those aspects are.

In this work (in progress), we apply this approach---developed for
introductory programming in our companion paper
\autocite{VTProgMisconceptions}---to the version control system Git: the
commit-graph model, the staging area, branches, and remotes. Git is a
grateful subject: its learning difficulties are documented (students
conceptualise branches as directories and try to \texttt{cd} into them
\autocite{IsomottonenCochez2014}), and part of the difficulty has been
located in the tool's own conceptual design rather than in the learner
\autocite{DeRossoJackson2016Gitless}. We treat the documented
misconceptions as phenomenographic observations, analyse them through the
lens of variation theory to identify critical aspects, and outline patterns
of variation to teach learners to discern them.
